\documentclass[./../main.tex]{subfiles} % 继承主文件设置
\begin{document}

\section{chapter1: 导论}

\begin{theorem}
    Stirling 公式：
    $$
    n! \sim \sqrt{2\pi n}\left(\frac{n}{e}\right)^n(1+\Theta(\frac{1}{n}))
    $$
\end{theorem}

\subsection{基本概念}
算法是在有限的步骤内求解某一问题，所使用的一组定义明确的规则。也就是在通过一组有限定义的指令来完成一些任务，给定一个初始状态，并相应的获取最终结果。

算法应该具有以下特性：可行性、确定性（每一步都被精确定义，不存在模糊操作）、有限性（程序可以没有这个特性，操作系统就是一个不关机就无限期执行的程序）、输入、输出。

算法是解决问题的抽象步骤，程序是算法的具体实现，算法通过编程语言转为程序，程序在计算机上运行。

算法分析有事前分析与事后测试两种基本方法。

\begin{definition}
    对于任何一个计算步骤，如果它的一次执行代价总是以一个时间常数为上界，与输入数据或者使用何种算法无关，则称该计算步骤为“元运算”或“元操作”。
\end{definition}
\begin{definition}
    算法中某个语句或运算的执行次数称为其频度。
\end{definition}

\begin{definition}
    如果算法中的一个元运算具有最高频度，所有其他元运算频度均在它的频度的常数倍内，则称这个元运算为基本运算。
\end{definition}

\subsection{复杂度}
\subsubsection{$O$记号：渐进上界记号}
设$f(n)$和$g(n)$是从自然数到非负正数集上的两个函数。如果存在正常数$c$和自然数$n_0$，使得
$$
\forall n\ge n_0, f(n)\le c\cdot g(n)
$$
则称$f(n)$是$O(g(n))$的，记为$f(n)=O(g(n))$\footnote{这种"="的记号的本意其实是$f(n)\in O(g(n))$，而对于出现在等式或不等式中的渐进记号表示的也是，其中的某一函数}。也可以这样理解，当$\lim_{n\to\infty}\frac{f(n)}{g(n)}$存在且不为无穷，那么就有$f(n)=O(g(n))$，(但是ppt上似乎还要求极限不为无穷，私以为这里把极限为无穷也算作了极限存在的一种)

\subsubsection{$\Omega$记号：渐进下界记号}

设$f(n)$和$g(n)$是从自然数到非负正数集上的两个函数，如果存在正常数$c$和自然数$n_0$，使得
$$
\forall ,\ge n_0, f(n)\ge \cdot g(n)
$$
则称$f(n)$是$\Omega(g(n))$的，记为$f(n)=\Omega(g(n))$。也可以这样理解，当$\lim_{n\to \infty}\frac{f(n)}{g(n)}$存在且不为0，则表示$f(n)=\Omega(g(n))$


\subsubsection{$\Theta$-记号：紧渐进界记号}
如果存在2个正常数$c_1,c_2$和自然数$n_0$，使得
$$
\forall n \ge n_0, c_1\cdot g(n)\le f(n)\le c_2\cdot g(n)
$$
则称$f(n)$是$\Theta(g(n))$，记为$f(n)=\Theta(g(n))$，也就是说，$\lim_{n\to\infty}\frac{f(n)}{g(n)}=c>0$, 则$f(n)=\Theta(g(n))$

可以推断，$f(n)=\Theta(g(n))\iff f(n)=O(g(n))\land f(n)=\Omega(g(n))$, 此时称$f(n)$与$g(n)$同阶，

\subsubsection{$o$-记号}
对于任意给定的$c>0$，都存在正整数$n_0$，使得当$n\ge n_0$时有
$$
f(n)<c\cdot g(n)
$$
则称函数$f(n)$是$o(g(n))$的，记为$f(n)=o(g(n))$也就是说$\lim_{n\to\infty}\frac{f(n)}{g(n)}=0$，则$f(n)=o(g(n))$.

因此可以有：
$$
f(n)=o(g(n))\iff f(n)=O(g(n))\land g(n)\neq O(f(n))
$$
当$n$充分大时，$f(n)$的阶比$g(n)$低。

\subsubsection{$\omega$-记号：非紧下界记号}
对于任意给定的$c>0$，都存在正整数$n_0$，使得当$n\ge n_0$时有$f(n)>\cdot g(n)$，则称函数$f(n)$时$\omega(g(n))$的，记为$f(n)=\omega(g(n))$,当$\lim_{n\to\infty}\frac{f(n)}{g(n)}$不存在，$\lim_{n\to\infty}\frac{f(n)}{g(n)}\to\infty, f(n)=\omega(g(n))$.

\subsubsection{渐进记号的性质}
\begin{enumerate}
    \item 传递性：所有符号
    \item 反身性：$\Theta, O, \Omega$
    \item 对称性：$\Theta$
    \item 互对称性：
    \begin{align*}
        f(n)=O(g(n))&\iff g(n)=\Omega(f(n))\\
        f(n)=o(g(n))&\iff g(n)=\omega(f(n))
    \end{align*}
    \item 算数计算
    \begin{itemize}
        \item $O(f(n))+O(g(n))=O(\max\{f(n),g(n)\})$
        \item $O(f(n))+O(g(n))=O(f(n)+g(n))$
        \item $O(f(n))*O(g(n))=O(f(n)*g(n))$
        \item $O(c\cdot f(n))=O(f(n))$
        \item $g(n)=O(f(n))\Rightarrow O(f(n))+O(g(n))=O(f(n))$
    \end{itemize}
\end{enumerate}

算法的空间复杂度是对一个算法在运行过程中\textbf{临时使用存储空间的渐进度量}

\subsection{平摊分析}
\subsubsection{聚集分析}
将所有的操作加起来，然后除以操作数，得到的结果就是平摊代价。

\subsubsection{记账法}
对不同操作赋予不同的费用，这个均摊成本$\hat c_{op}$可能与实际成本$c_{op}$不同，但必须要满足的是，对于任意的合法的$n$个操作序列，必须有$T(n)=\sum_{i=1}^nc_i\le \sum_{i=1}^n\hat c_i$.

这种方法更常见于对一个已知的均摊代价进行巧妙的证明。例如：对于二进制计数器算法，因为每一个操作1翻转为0中的1，都一定来自于一个操作：0翻转为1，而最终结果一定会剩余一定量的1，所以操作：0翻1的数量大于1翻0的数量。

\begin{table}[h]
\centering
\caption{Binary Counter Flip Operations: Real Cost vs. Amortized Cost}
\begin{tabular}{lcc}
\toprule
\textbf{OP} & \textbf{Real Cost }$\boldsymbol{C_{OP}}$ & \textbf{Amortized Cost }$\boldsymbol{\widehat{C}_{OP}}$ \\
\midrule
flip ($0 \to 1$) & 1 & 2 \\
flip ($1 \to 0$) & 1 & 0 \\
\bottomrule
\end{tabular}
\end{table}

\begin{align*}
T(n) &= \sum_{i=1}^{n} C_i \\
     &= \#\mathrm{flip}(0 \rightarrow 1) + \#\mathrm{flip}(1 \rightarrow 0) \\
     &\leq 2 \cdot \#\mathrm{flip}(0 \rightarrow 1) \\
     &\leq 2n
\end{align*}


\subsubsection{势能法}
势能函数充当能量存储罐的作用，当操作简单时就会势能增加（相当于存钱）；当操作复杂时就会势能减少（相当于取钱）；那么每次操作的平摊代价=实际操作成本+存储罐的变化量。即：
$$
a_i=c_i+\Phi(D_i)-\Phi(D_{i-1})
$$
那么对于$n$个操作的总平摊代价为：
$$
\sum a_i=\sum c_i+\Phi(D_n)-\Phi(D_0)
$$

如果势函数使得对所有的$n$，都有$\Phi(D_n)\ge\Phi(D_0)$，则总平摊代价就是总实际代价的一个上界，为了方便起见有时会定义$\Phi(D_0)=0$.

\subsection{递推关系}
（高中部分内容不再展开）

\subsubsection{主定理}
\begin{theorem}
    设$a\ge1,b\ge1$为常数，$f(n)$为函数, $T(n)$为非负正数，且
    $$
    T(n)=aT(\frac{n}{b})+f(n)
    $$
    那么就有
    \begin{itemize}
        \item $f(n)=O(n^{\log_ba-\varepsilon}), \varepsilon>0$
        那么$T(n)=\Theta(n^{\log_ba})$
        \item $f(n)=\Theta(n^{\log_ba})$,
        那么$T(n)=\Theta(n^{\log_ba}\log n)$
        \item $f(n)=\Omega(n^{\log_ba+\varepsilon}), \varepsilon>0$，
        且对某个常数$c<1$和所有充分大的$n$有$a\cdot f(\frac{n}{b})\le c\cdot f(n)$，
        那么$T(n)=\Theta(f(n))$
    \end{itemize}
\end{theorem}
\begin{proof}
    设$n=b^k$, 则有
    \begin{align*}
        T(n)&=a^kT(1)+\sum_{j=0}^{k-1}a^jf(\frac{n}{b^j})\\
        &=c_1n^{\log_ba}+\sum_{j=0}^{k-1}a^jf(\frac{n}{b^j}), T(1)=c_1\\
    \end{align*}

    \begin{itemize}
        \item \textbf{情况1}
        \begin{align*}
            T(n) &= c_{1}n^{\log_b a} + \sum_{j=0}^{k-1} a^{j}f\left(\frac{n}{b^{j}}\right) \\
            &= c_{1}n^{\log_b a} + O\left(\sum_{j=0}^{\log_b n -1} a^{j}\left(\frac{n}{b^{j}}\right)^{\log_b a - \varepsilon}\right) \\
            &= c_{1}n^{\log_b a} + O\left(n^{\log_b a - \varepsilon} \sum_{j=0}^{\log_b n -1} \frac{a^{j}}{(b^{\log_b a - \varepsilon})^{j}}\right) \\
            &= c_{1}n^{\log_b a} + O\left(n^{\log_b a - \varepsilon} \sum_{j=0}^{\log_b n -1} (b^{\varepsilon})^{j}\right) \\
            &= c_{1}n^{\log_b a} + O\left(n^{\log_b a - \varepsilon} \cdot \frac{b^{\varepsilon \log_b n} - 1}{b^{\varepsilon} - 1}\right) \\
            &= c_{1}n^{\log_b a} + O\left(n^{\log_b a - \varepsilon} \cdot n^{\varepsilon}\right) \\
            &= \Theta(n^{\log_b a})
        \end{align*}
        \item \textbf{情况2}

        \begin{align*}
            T(n) &= c_{1}n^{\log_{b}a} + \sum_{j=0}^{k-1} a^{j}f\left(\frac{n}{b^{j}}\right) \\
            &= c_{1}n^{\log_{b}a} + \Theta\left(\sum_{j=0}^{\log_{b}n-1} a^{j}\left(\frac{n}{b^{j}}\right)^{\log_{b}a}\right) \\
            &= c_{1}n^{\log_{b}a} + \Theta\left(n^{\log_{b}a}\sum_{j=0}^{\log_{b}n-1} \frac{a^{j}}{a^{j}}\right) \\
            &= c_{1}n^{\log_{b}a} + \Theta\left(n^{\log_{b}a}\log n\right) \\
            &= \Theta\left(n^{\log_{b}a}\log n\right)
        \end{align*}
        \item \textbf{情况3}
        \begin{align*}
            T(n) &= c_{1}n^{\log_{b}a} + \sum_{j=0}^{k-1} a^{j}f\left(\frac{n}{b^{j}}\right) \\
            &\leq c_{1}n^{\log_{b}a} + \sum_{j=0}^{\log_{b}n-1}c^{j}f(n) \quad \left(a f\left(\frac{n}{b}\right)\leq cf(n)\right) \\
            &= c_{1}n^{\log_{b}a} + f(n)\frac{c^{\log_{b}n}-1}{c-1} \\
            &= c_{1}n^{\log_{b}a} + \Theta(f(n)) \quad (c < 1) \\
            &= \Theta(f(n)) \quad \left(f(n)=\Omega\left(n^{\log_{b}a+\varepsilon}\right)\right)
        \end{align*}
    \end{itemize}
\end{proof}

对于$T(n)=2T(\frac{n}{2})+n\log n$, 就是不能使用主定理的，因为不存在$\varepsilon>0$使得$n\log n=\Omega(n^{1+\varepsilon})$.

这个可以用递归树进行求解，恰好这个拆解得到的是平衡树，只需要对$f(n)$项求和即可。
\begin{align*}
    T(n)&=\sum_{i=0}^{\log n}n(\log n-i)\\
    &=n\log^2n=O(n\log^2n)
\end{align*}

\subsubsection{Akra-Bazzi定理}
Akra-Bazzi定理是分析分治算法时间复杂度的强大工具，它推广了经典的主定理(Master Theorem)，适用于更广泛的递归关系。

对于递归关系式：
$$ T(n) = \sum_{i=1}^{k} a_iT(b_in + h_i(n)) + f(n) $$
其中：
\begin{itemize}
\item $a_i > 0$ 且 $0 < b_i < 1$（对所有$i=1,\ldots,k$）
\item $|h_i(n)| = O\left(\frac{n}{\log^2 n}\right)$（扰动项）
\item $f(n)$ 满足某些正则条件
\end{itemize}
定义$p$为方程的解：
$$ \sum_{i=1}^{k} a_ib_i^p = 1 $$
则时间复杂度为：
$$ T(n) = \Theta\left(n^p\left(1 + \int_1^n \frac{f(u)}{u^{p+1}}du\right)\right) $$

\begin{table}[h]
\centering
\caption{主定理与Akra-Bazzi定理对比}
\begin{tabular}{|l|l|}
\hline
\textbf{特征} & \textbf{Akra-Bazzi定理} \\ \hline
适用范围 & 非均匀划分的递归式 \\ \hline
扰动项 & 允许$h_i(n)$存在 \\ \hline
驱动函数 & $f(n)$形式更灵活 \\ \hline
解的形式 & 积分表达式 \\ \hline
\end{tabular}
\end{table}


\subsubsection{应用示例}
分析递归式 $T(n) = T(\lfloor n/2 \rfloor) + T(\lceil n/2 \rceil) + n$：
\begin{enumerate}
\item 确定参数：$a_1=a_2=1$, $b_1=b_2=1/2$
\item 解方程：$2\cdot(1/2)^p = 1 \Rightarrow p=1$
\item 计算积分：$\int_1^n \frac{u}{u^2}du = \ln n$
\item 最终解：$T(n) = \Theta(n(1 + \ln n)) = \Theta(n\log n)$
\end{enumerate}

当递归式满足主定理条件时，Akra-Bazzi定理给出的结果与主定理一致，但前者适用性更广。

不过其实用数学归纳法会更舒适一些。例如估计以下复杂度：
$$
T(n)=T(\frac{n}{5}) + T(\frac{7n}{10}) + \Theta(n)
$$

假设对任意的$m<n$，都有$T(m)<cm$恒成立，其中$c$为常数，下面证明$T(n)<cn$即可证明出$T(n)=\Theta(n)$

$$
\begin{aligned}
    T(n)&=T(\frac{n}{5}) + T(\frac{7n}{10}) + \Theta(n)\\
    &<T(\frac{n}{5}) + T(\frac{7n}{10}) + dn\\
    &<c\frac{n}{5} + c\frac{7n}{10} + dn\\
    &=(\frac{c}{5}+\frac{7c}{10}+d)n<cn\\
\end{aligned}
$$

因而只需取$c>10d$即可。

\subsubsection{递推关系求解练习}
\begin{enumerate}
    \item $T(n)=T(\frac{n}2)+T(\frac{n}4)+cn$
    \begin{align*}
        &\text{令}n=2^k, T(2^k)=T(2^{k-1})+T(2^{k-2})+c\cdot 2^k\\
        &\text{构建递归树知，若设}T(2^{k-i})\text{的结点数量是}f[i]\text{，则有}f[i]=f[i-1]+f[i-2]\\
        &\text{且}f[0]=f[1]=1, \text{因此叶子结点数量为}f[k]=\frac{1}{\sqrt{5}}\left[\left(\frac{1+\sqrt{5}}{2}\right)^{k+1}-\left(\frac{1-\sqrt{5}}{2}\right)^{k+1}\right]\\
        &\frac{1}{\sqrt{5}}\left[\left(\frac{1+\sqrt{5}}{2}\right)^{k+1}-\left|\left(\frac{1-\sqrt{5}}{2}\right)^{k+1}\right|\right]\leq f[k]\leq\frac{1}{\sqrt{5}}\left(\frac{1+\sqrt{5}}{2}\right)^{k+1}\\
        &\therefore f[k]=O(a^k)=O(n^{\log_2 a}), \text{其中}a=\frac{1+\sqrt{5}}{2}\\
        &\text{而对于计算耗时为：}cn\cdot\sum_{i=1}^n\left(\frac{3}{4}\right)^{i-1}=O(n)>O(n^{\log_2 a})\\
        &\therefore T(n)=O(n)
    \end{align*}
    \item $T(n)=T(n-1)+\log n$
    \begin{align*}
        &\text{原式}=T(1)+\sum_{i=2}^n\log i\\
        \text{下证}&: \text{$\sum_{i=2}^n \log i$ 与 $n \log n$ 同阶} \\
        \text{由Stolz定理知，}& \lim_{n\to\infty} \frac{\sum_{i=2}^n \log i}{n \log n} = \lim_{n\to\infty} \frac{\log n}{n \log n - (n-1)\log(n-1)} = \lim_{n\to\infty} \frac{\log n}{\log(n-1)} = 1 \\
        &\therefore T(n) = \Theta(n \log n)
    \end{align*}
\end{enumerate}

\subsection{本章作业}

\subsubsection{势函数规范化构造及其平摊代价不变性证明}
假定有势函数 $\Phi$，对所有 $i \geq 1$ 满足 $\Phi(D_{i}) \geq \Phi(D_{0})$，但 $\Phi(D_{0}) \neq 0$。证明：存在势函数 $\Phi'$，使得 $\Phi'(D_{0}) = 0$，对于所有 $i \geq 1$，满足 $\Phi'(D_{i}) \geq 0$，且使用 $\Phi'$ 的平摊代价与使用 $\Phi$ 的平摊代价相同。

\begin{proof}
设新的势函数为：
$$
\Phi'(D_i) = \Phi(D_i) - \Phi(D_0) \quad (i = 0,1,2,\ldots)
$$

\begin{enumerate}
\item \textbf{零势初始化}：
$$
\Phi'(D_0) = \Phi(D_0) - \Phi(D_0) = 0
$$

\item \textbf{非负性保持}（$\forall i \geq 1$）：
$$
\Phi'(D_i) = \Phi(D_i) - \Phi(D_0) \geq 0 \quad \text{（由已知条件 $\Phi(D_i) \geq \Phi(D_0)$）}
$$

\item \textbf{平摊代价不变性}：
对于任意操作，设实际代价为 $c_i$，有：
$$
\begin{aligned}
\hat{c}_i' &= c_i + \Phi'(D_i) - \Phi'(D_{i-1}) \\
&= c_i + (\Phi(D_i) - \Phi(D_0)) - (\Phi(D_{i-1}) - \Phi(D_0)) \\
&= c_i + \Phi(D_i) - \Phi(D_{i-1}) \\
&= \hat{c}_i
\end{aligned}
$$
\end{enumerate}
\end{proof}

\subsubsection{聚集法计算平摊代价}
假定对一个数据结构执行一个由$n$个操作组成的操作序列，当$i$严格为2的幂时，第$i$个操作的代价为$i$，否则代价为1。请使用聚集法确定每个操作的平摊代价。

考虑$n$个操作序列的总代价：
$$
T(n) = \sum_{\substack{k=0 \\ 2^k \leq n}}^{\lfloor \log_2 n \rfloor} 2^k + \Big(n - \lfloor \log_2 n \rfloor -1\Big)
$$
展开后可得：
$$
T(n) = (2^{\lfloor \log_2 n \rfloor +1} -1) + (n - \lfloor \log_2 n \rfloor -1) \leq 3n
$$
因此平摊代价为：
$$
\frac{T(n)}{n} \leq 3 = O(1)
$$

\subsubsection{主定理运用}
设递推方程 $T(n)=7T(\frac{n}{2})+n^{2}$ 给出了算法$A$在最坏情况下的时间复杂度函数，算法$B$在最坏情况下的时间复杂度函数 $W(n)$ 满足递推方程 $W(n)=aW(\frac{n}{4})+n^{2}$。试确定最大的正整数 $a$，使得 $W(n)$ 的阶低于 $T(n)$ 的阶。

\begin{align*}
    &\because T(n) = 7T\left(\frac{n}{2}\right) + n^{2} \\
    &\therefore a = 7,\ b = 2,\ f(n) = n^{2} \\
    &\therefore n^{\log_{b}a} = n^{\log_{2}7},\ 
    f(n) = n^{2} = O\left(n^{\log_{2}7 - \varepsilon}\right),\ \varepsilon > 0 \\
    &\therefore T(n) = \Theta\left(n^{\log_{2}7}\right) \\
    &\text{又}\because W(n) = aW\left(\frac{n}{4}\right) + n^{2} \\
    &\therefore n^{\log_{b}a} = n^{\log_{4}a} = n^{\frac{1}{2}\log_{2}a},\ 
    f'(n) = n^{2} \\
    &\text{if } a = 16,\ f'(n) = \Theta\left(n^{\log_{b}a}\right),\ 
    W(n) = \Theta\left(n^{2} \log n\right) < \Theta\left(n^{\log_{2}7}\right),\ \text{满足题意} \\
    &\text{if } a > 16,\ f'(n) = O\left(n^{\log_{b}a - \varepsilon}\right),\ \varepsilon > 0,\ 
    W(n) = \Theta\left(n^{\log_{b}a}\right) < \Theta\left(n^{\log_{2}7}\right) \\
    &\therefore \sqrt{a} < 7,\ a < 49 \\
    &\therefore a_{\max} = 48 > 16 \\
    &\therefore a_{\max} = 48
\end{align*}
注意$a=16$的讨论是必要的。


\end{document}